% ============================================================
% Chapter 4: Proposed Method: Specification and Design
%   Section: Interface Design
% Included from chapters/04-proposed-method/proposed-method.tex
% Source: frontend/CLAUDE.md.
% ============================================================

\section{Interface Design}
\label{sec:design-interface}

The interface serves three audiences, and its pages are grouped by them.

\paragraph{The newsroom.} Four pages follow the editor's work in order,
with a step bar on each:

\begin{enumerate}
    \item \textbf{Discover}: pick one to three topic groups, set the
    balance between positive impact and topic relevance, and read the
    ranked candidates of a round, each with its place, score and the
    parts of its score. The AI selection can be asked for; its picks
    arrive ticked with a score and a reason, and stay editable. At most
    twenty are sent to analysis. The round is in the page's address, so
    a reload resumes it.
    \item \textbf{Analyze}: paste article URLs and follow each analysis
    as it happens; claims appear as rows the moment they are selected
    and fill in with their own verdicts as they finish.
    \item \textbf{Live}: a second screen that follows any run from any
    client --- the queries sent, every source found and the engines
    behind it, the rating each received, and the verdict.
    \item \textbf{Reader}: what a reader would see.
\end{enumerate}

\paragraph{The lab.} One page per stage, for inspecting and debugging
it in isolation: a single claim to a verdict with the model's answer
before recalibration; the enrichment of a URL or a text; the text
corrector; the scraper's request statistics, source health and time to
reception; a health check of every source; the graph, with its schema,
a read-only query console and related articles; and the evaluation's
state --- the labelled set against its targets, the claim selector's
precision and the harness's runs with their intervals.

\paragraph{The reader.} The public reader view lists the articles the
pipeline marked publishable, filtered by topic and language. An article
shows each claim's verdict and confidence, the model's reasoning, and
its sources with their stance, the quoted passage, the site's
reliability (or ``unrated'') and why each was found. A claim whose
search or model failed reads ``not checked'' and is counted apart from
the verdicts. Only the article's lead and a link are shown, never its
body or a scraped page (Section~\ref{subsec:third-party-content}).

The design follows two rules. Uncertainty is shown, not hidden: every
score carries its parts, every rating says whether it is real or a
default, every outage says so. And the interface is a client of the
same API the evaluation uses, so what a person sees is what is measured.
